\begin{abstract}
Multilingual users often interact with large language models (LLMs) through
code-switched prompts, mixed-language editing requests, and constraints about
script, register, locale, and literal text preservation. Standard one-turn
evaluations usually ask whether the final answer is correct, but they miss a
common usability failure: the model returns a fluent, plausible answer that
violates the user's interaction contract, forcing the user to spend extra turns
repairing the response. We introduce \emph{Re-prompt Tax}, a metric family that
measures first-turn global alignment, the number of repair turns needed for
recovery, unresolved failures, and token overhead. We construct
\textsc{RePromptTax-Stress-v0.2}, a 120-item synthetic stress pilot across
Spanish-English, Hindi-English, and Arabic-English interactions, and evaluate
three GPT-4.1-family API models under a baseline prompt and a Global Interaction
Contract prompt. Baseline first-turn global alignment ranges from 67.5\% to
76.7\%; the mitigation improves alignment to 76.7--93.3\% and reduces mean
repair turns and repair-normalized token tax. A current-model refresh shows the
same pattern on full 120-item \texttt{gpt-5.4-mini} and \texttt{gpt-5.5} runs:
under \texttt{gpt-5.5}, alignment rises from 81.7\% to 98.3\%. Prompt-mechanism
diagnostics show that a narrower content-preservation scaffold captures most of
this gain, but only near-ties the full contract on current models. Analysis
shows that the dominant failure mode is not generic multilingual QA, but
implicit editing preservation: models often translate or frame content in the
instruction language when the user intended the quoted content to remain in
English. A 10\% stratified blinded
LLM-judge audit agrees with the automatic scorer on 71/72 sampled pass/fail
labels, and a \texttt{gpt-5.5} judge refresh agrees on 70/72. We release a
conservative evaluation protocol and claim boundary for
measuring hidden global usability costs rather than broad multilingual ability.
\end{abstract}
